\section{Broadcasting a shared label}\label{sec:copy}

Suppose that every source in \(P\) knows the same \(b\)-bit label \(x\).
To send \(x\) to a set \(Z\) of new sites, one source can select the
\(x\)-term of a diagonal shared by \(P\cup Z\). This costs \(w(P\cup Z)b\).
Agreement across the sources lets us recover a diagonal on \(P\) in the
same operation. The returned diagonal has price \(w(P)b\).
This use of a readable label is related to
\cite[Proposition~\ext{G}{lem:submod}]{oaG}.

\begin{lemma}[Broadcasting]\label{lem:copy}
Fix a spectral point on \(S\), disjoint nonempty sets \(P,Z\subseteq S\),
and an integer \(b\ge1\). Put \(K=P\cup Z\), \(N=2^b\), and
\(\ell=\log_2(2|K|)\). Give each site of \(P\) a control in \([N]\).
There are an input tensor \(R\) and controlled local maps that return
\(G_P(N)\otimes E_Z(x)\) whenever the controls all equal \(x\), with
\[
 h(R)-h(G_P(N))
       \le (w(K)-w(P))b+O\!\left(w(K)\sqrt{b\ell}\right).
\]
The implied constant is absolute. Any such conversion returning
\(Q\otimes E_Z(x)\) with \(Q\) fixed and nonzero has net cost at least
\((w(K)-w(P))b\).
\end{lemma}

\begin{proof}
For the lower bound, supply the common control as the coherent sum
\(G_P(N)\) and retain its local basis vectors. Term by term, the conversion
extends the common label to \(Z\). Thus
\(G_P(N)\otimes R\longrightarrow G_K(N)\otimes Q\), and monotonicity
gives the bound.

For the construction, write \(a=|P|\) and \(c=|Z|\). The elementary
broadcast described above selects \(x\) at one source and sends every
basis vector to \(1\) at the others. Carry an independent \(G_P(N)\)
through this operation. Its net cost \(w(K)b\) suffices when \(a=1\),
since then \(w(P)=0\). It also suffices when \(b\le\ell\), since then
\(w(P)b\le w(K)\sqrt{b\ell}\). Henceforth \(a\ge2\) and \(b>\ell\).

We encode the broadcast label and the returned diagonal by integer
vectors. In a box \(V=\{0,\ldots,q-1\}^m\), choose a set \(H\) of
\(N\) vectors with the same Euclidean norm and identify \([N]\) with \(H\).
For each \(u\in H\), choose an ordered set \(V_u\subseteq V\) of \(N\)
vectors having the same norm and the same inner product with \(u\).
We choose \(q,m\) below. On control \(u\), source \(j\) will output
the index of a vector \(v_j\in V_u\), and target \(z\) will output
the label \(g_z\in H\). Impose the single vector equation
\begin{equation}\label{eq:broadcast-relation}
  \sum_{z\in Z}g_z=cu+\sum_{j=2}^a v_j-(a-1)v_1 .
\end{equation}

Taking the inner product with \(u\) cancels all the terms involving
\(v_j\). It gives \(\sum_z\langle u,g_z\rangle=c\|u\|^2\).
Since every \(g_z\) has the norm of \(u\), equality forces \(g_z=u\)
for every \(z\). Equation~\eqref{eq:broadcast-relation} now says that
\(v_1\) is the average of the other source vectors. Their equal norms
give \(\sum_{j=2}^a\|v_j-v_1\|^2=0\). Conversely, every
\(v\in V_u\) gives a solution with all \(v_j=v\) and all \(g_z=u\).
Summing over the solutions therefore returns exactly \(N\) diagonal
terms on \(P\), each with the desired label at every site of \(Z\).
The use of a sphere to force equality in an average is the geometric
step in Behrend's construction~\cite{Behrend}.

To impose the equation by local maps, take \(L=(a+c-1)(q-1)+1\),
increasing it by one if the characteristic divides it. Write
\(\Gamma=(\mathbb Z/L\mathbb Z)^m\). In every coordinate the two
sides of Equation~\eqref{eq:broadcast-relation} differ by less than
\(L\) in absolute value. Their equality can therefore be tested in
\(\Gamma\). For \(D\in\Gamma\) the indicator of \(D=0\) is
\(|\Gamma|^{-1}\sum_{\chi\in\widehat\Gamma}\chi(D)\), since the
field contains all \(L\)-th roots of unity and \(L\) is invertible.

Supply \(R=G_K(|\Gamma|)\) with its shared index named by \(\chi\).
Each target sums over its possible \(g_z\) with coefficient \(\chi(g_z)\).
Each source \(j\ge2\) sums over \(v_j\in V_u\) with coefficient
\(\chi(-v_j)\), and source \(1\) uses
\(|\Gamma|^{-1}\chi((a-1)v_1-cu)\).
The product of these coefficients is the character summand for the
difference of the two sides. Each output index is the basis label
specified above, so character orthogonality gives
\(G_P(N)\otimes E_Z(u)\). Only the source maps use a control.

It remains to choose a small enough group. Set
\[
 d=\lceil\sqrt{b\ell}\rceil+4,\qquad q=2^d,\qquad
 m=\lceil b/d\rceil+6 .
\]
The squared norm of a point in \(V\) and its inner product with a
fixed \(u\in V\) each take at most \(B=1+m(q-1)^2\le q^3\) values,
since \(m\le d+6\le q\). Pigeonholing by norm gives \(H\), and
pigeonholing by both statistics gives every \(V_u\), because
\(|V|/B^2\ge q^{m-6}\ge 2^b\). Finally \(\log_2 L\le d+\ell\) and
\[
 \log_2|\Gamma|\le (b/d+7)(d+\ell)
                  \le b+O(\sqrt{b\ell}).
\]
Subtracting the price \(w(P)b\) of the returned diagonal gives the
claimed net cost.
\end{proof}

\section{Constructing the equality test}\label{sec:filter}

On the promise of Theorem~\ref{thm:filter}, let \(x\) be the label
shared by \(P\) and \(y\) the label shared by \(T\). A perfect matching
on the alphabet contracts each pair to a single name, cutting the
number of bits by one. We will run the shorter equality test on a
family of such matchings. When \(x\ne y\), exactly one matching
pairs them. Its index can be computed by two broadcasts, so the
resulting false match can be cancelled exactly.

\begin{proof}[Proof of Theorem~\ref{thm:filter}]
Put \(K=P\cup T\) and \(g=w(P)+w(T)-w(K)\).
Fixed equal controls take \(R\) to \(Q\), so every filter has
nonnegative cost. Supply the two controls independently as
\(G_P(2^b)\otimes G_T(2^b)\), retaining their local indices.
The filter kills all terms with different labels, giving
\[
 G_P(2^b)\otimes G_T(2^b)\otimes R
       \longrightarrow G_K(2^b)\otimes Q .
\]
Taking prices gives the second lower bound \(h(R)-h(Q)\ge -gb\).

We construct the upper bound first for \(S=K\), so \(w(K)=1\).
Fix a bit length \(b\ge1\). The labels form the set \(\mathbb F_2^b\),
regardless of the characteristic of the tensor field. For each
nonzero \(p\), the pairs \(\{x,x+p\}\) form a perfect matching.
Name its pairs by \((b-1)\)-bit strings and write \(\pi_p(x)\)
for the name of the pair containing \(x\). If \(x=y\), the names
agree in every matching. If \(x\ne y\), they agree precisely when
\(p=x+y\). Keeping the index of each surviving matching gives the
identity
\[
 \sum_{\substack{p\ne0\\ \pi_p(x)=\pi_p(y)}} E_K(p)
 =
 \begin{cases}
 G_K(2^b-1), & x=y,\\
 E_K(x+y),   & x\ne y .
 \end{cases}
\]

We need a correction producing the second line on unequal inputs and
zero on equal inputs. Broadcast \(x\) from \(P\) to \(T\) by
Lemma~\ref{lem:copy}. At each site of \(T\) the control \(y\)
selects the map \(e_x\mapsto e_{x+y}\) for \(x\ne y\) and
\(e_y\mapsto0\). Keep this new basis label at \(T\) and use it as
the common control for a second broadcast from \(T\) to \(P\).
Denote the tensor products of the two input resources and the two
returns by \(U_b\) and \(V_b\), respectively. On unequal inputs
the output is \(V_b\otimes E_K(x+y)\). All its label slots may be
projected onto the span of the nonzero indices. With
\(\ell=\log_2(2|K|)\), Lemma~\ref{lem:copy} gives
\begin{equation}\label{eq:filter-correction}
 h(U_b)-h(V_b)\le (1-g)b+C_0\sqrt{b\ell}.
\end{equation}
Indeed, the two leading costs add to
\((1-w(P))b+(1-w(T))b=(1-g)b\).

Start the recursion with \(R_0=Q_0=1\). At bit length \(b\), let
\(L_b=2^b-1\) and give every site access to the same matching index
through \(G_K(L_b)\). The shorter filter on the two pair names
can be applied under control of this index. With \(p\) retained,
its output on equal labels contains the entire \(G_K(L_b)\).
Including that diagonal in the return leads to the following choice
of resources for the two branches.
\begin{align*}
 Q_b &= Q_{b-1}\otimes V_b\otimes G_K(L_b),\\
 R_b &= \bigl(R_{b-1}\otimes V_b\otimes G_K(L_b)\bigr)
             \oplus \bigl(Q_{b-1}\otimes U_b\bigr).
\end{align*}
The first summand runs the shorter filter, keeps \(V_b\), and
returns \(Q_b\) on equal controls. On unequal controls its output is
\(Q_{b-1}\otimes V_b\otimes E_K(x+y)\). The second summand retains
\(Q_{b-1}\) and forms exactly this correction, with zero output on
equal controls. Multiply its map at one site by \(-1\).
Both sets of local maps now land in the same spaces, so their maps on
the local direct sums produce the difference. This eliminates the
single false match in every characteristic.

Set \(f_b=\varphi(R_b)/\varphi(Q_b)\). The diagonal of \(L_b\)
matchings is furnished and returned by the recursive branch. Thus
only the correction contributes a new term to the price ratio.
Additivity on direct sums and Equation~\eqref{eq:filter-correction}
give
\[
 f_0=1,\qquad
 f_b=f_{b-1}+\frac{\varphi(U_b)}{L_b\varphi(V_b)}
      \le f_{b-1}+2^{1-gb+C_0\sqrt{b\ell}} .
\]
Here \(L_b\ge 2^{b-1}\). Each of the first \(b\) increments is at most
\(2^{1+\max\{0,-g\}b+C_0\sqrt{b\ell}}\). Taking logarithms of their
sum with \(f_0\) and using \(\log_2(b+1)=O(\sqrt b)\) yields
\[
 \log_2 f_b\le\max\{0,-g\}b+C_1\sqrt{b\ell}
\]
with absolute \(C_1\). When \(g>0\), the entire series is bounded.
Indeed, \(C_0\sqrt{b\ell}\le gb/2+C_0^2\ell/(2g)\) gives
\[
 f_b\le 1+\frac{2^{1+C_0^2\ell/(2g)}}{2^{g/2}-1}.
\]
Since \(0<g\le1\) and \(\ell\ge1\), this gives
\(\log_2 f_b=O(\ell/g)\) with an absolute constant.

To place this filter in the original set of sites \(S\), put
\(\theta=w(K)\). If \(\theta=0\), the tensor \(G_K(2^b)\) has price
zero. One site in each group selects its control index from that
diagonal and the other sites send every index to \(1\), giving
\(\mathbf 1_{\{x=y\}}\) at zero cost.

Suppose \(\theta>0\). Write \(\iota X\) for a \(K\)-tensor padded by
scalar factors on \(S\setminus K\). Then
\(\psi(X)=\varphi(\iota X)^{1/\theta}\) is a spectral point on \(K\)
\cite[Section~\ext{G}{sec:spectra}]{oaG}. To see additivity, expand
\((X\oplus_K Y)^{\otimes m}\) and group terms by their number \(j\) of
\(X\)-factors. The \(j\)-th group is locally isomorphic to
\[
 G_K\!\left(\binom mj\right)
       \otimes X^{\otimes j}\otimes Y^{\otimes(m-j)}.
\]
Each group can be extracted by a projection. The direct sum on \(S\)
of all the padded groups restricts to the padded power by identifying
the scalar slots outside \(K\). Hence \(\varphi(\iota(X\oplus_K Y))^m\)
lies between the largest group value and \(m+1\) times that value.
By Equation~\eqref{eq:price}, the \(1/\theta\) power of the \(j\)-th value
is \(\binom mj\psi(X)^j\psi(Y)^{m-j}\). Its maximum differs from
the binomial sum by at most a factor \(m+1\). Taking \(m\)-th roots
proves additivity. The remaining spectral axioms follow from those
of \(\varphi\).

Under \(\psi\) the gap on \(K\) is \(g/\theta\). Padding the filter
constructed there multiplies its net cost by \(\theta\). Thus for
every \(b\ge1\) its cost on \(S\) exceeds \(\max\{0,-g\}b\) by at
most \(C_1 w(K)\sqrt{b\ell}\), and for \(g>0\) it is bounded by
\(O(w(K)^2\ell/g)\) uniformly in \(b\). The lower bounds proved
above then give the limiting cost per bit.
\end{proof}
